% ======================================================================
%  Supplementary Material / Appendix
%  VLM-Verified Counterfactual Hindsight for Sparse-Reward Manipulation
% ======================================================================
%
% Note: this appendix is intended to extend Appendix A in main.tex.
% Section A.1 (Prompt Template) and A.2 (Heuristic localizer) are already
% sketched in main.tex (\appendix block); we extend them here with full
% pseudocode, derivations, hyperparameter tables, and reproducibility
% material.

\section{Extended Method Details}
\label{app:methods}

\subsection{Algorithm pseudocode}
\label{app:algorithms}

We provide complete pseudocode for the two algorithmic contributions of
the paper: Semantic PER with the failure-direction signal
(Algorithm~\ref{alg:sem-per}) and VLM-Verified Counterfactual Hindsight
(Algorithm~\ref{alg:verified-cf}).

\begin{figure}[t]
\centering
\fbox{\parbox{0.95\linewidth}{
\textbf{Algorithm 1: Semantic PER with the failure-direction signal.}\\[2pt]
\textbf{Input.} Buffer $\mathcal{D}_B$ with sum-tree priorities $\{p_i\}$;
TD-error array $\{\delta_i\}$; semantic-weight array $\{w_i\}$ (init $1$);
SAC critic $Q_\theta$; PER exponent $\alpha$; IS exponent $\beta_t$; window
half-width $W$; cap $w_{\max}$; VLM $q_\phi$; localizer
$\mathrm{Loc}_\phi$.\\[2pt]
\textbf{1. Episode collection.} For step $t=0,\dots,T-1$: act
$a_t\!\sim\!\pi_\theta(\cdot\mid s_t)$, observe $(r_t,s_{t+1})$, push
$(s_t,a_t,r_t,s_{t+1})$ to $\mathcal{D}_B$ at index
$i_t\!\in\!\{0,\dots,|\mathcal{D}_B|\!-\!1\}$ with priority
$p_{i_t}=\max_j p_j$ (PER's max-priority initialization).\\[2pt]
\textbf{2. Failure check.} If episode ended without success ($r_T<0$): goto~3;
else skip to step~5.\\[2pt]
\textbf{3. VLM failure localization.} Select $K=5$ uniformly-spaced keyframes
$\{F_{k}\}_{k=0}^{K-1}$ from the episode RGB renders. Build
\texttt{frame\_index\_map}. Query
$t_{\text{fail}} \leftarrow \mathrm{Loc}_\phi(\{F_k\}, \text{task\_desc},
T)$ (returns keyframe index $\hat k$; map to timestep
$t_{\text{fail}}\leftarrow$ keyframe-timestep$(\hat k)$).\\[2pt]
\textbf{4. Apply semantic weight.} For each in-episode slot $i$ with episode-local
timestep $\tau_i\in\{0,\dots,T-1\}$:
$w_i \leftarrow w_{\max}$ if $|\tau_i - t_{\text{fail}}|\le W$ else $1$.
Update sum-tree leaf for slot $i$ to $p_i \leftarrow (|\delta_i|+\varepsilon)^\alpha
\cdot w_i$.\\[2pt]
\textbf{5. Critic update.} Sample minibatch $\{i_j\}_{j=1}^{B}$ from
$\mu(i)\propto p_i$. Compute IS weights
$w_{\text{IS},j} = (|\mathcal{D}_B|\cdot \mu(i_j))^{-\beta_t}$, normalize by
$\max_j w_{\text{IS},j}$. Compute TD-errors
$\delta_{i_j} = r_{i_j} + \gamma Q_{\bar\theta}(s_{i_j+1},
\pi_\theta(s_{i_j+1})) - Q_\theta(s_{i_j},a_{i_j})$, take loss
$\mathcal{L} = \frac{1}{B}\sum_j w_{\text{IS},j}\,\delta_{i_j}^2$.
Backprop; update $\theta$.\\[2pt]
\textbf{6. Priority update.} For each $j$: $\delta_{i_j}\!\to\!\text{stored TD}$,
set $p_{i_j} \leftarrow (|\delta_{i_j}|+\varepsilon)^\alpha \cdot w_{i_j}$.\\[2pt]
\textbf{7. Goto~1 until total\_steps reached.}\\[2pt]
\textbf{Output.} Trained policy $\pi_\theta$.
}}
\caption{Semantic PER with the failure-direction signal. The semantic-weight
multiplier $w_i$ is set once per episode at the conclusion of step~3 and
remains attached to slot $i$ until the slot is overwritten by the ring
buffer. Step~5's IS weights use the PER proposal $\mu_P$, not the full
$\mu_{\text{Sem}}$; see Section~\ref{sec:theory} and
Appendix~\ref{app:is-derivation} for the under-correction discussion.}
\label{alg:sem-per}
\end{figure}

\begin{figure}[t]
\centering
\fbox{\parbox{0.95\linewidth}{
\textbf{Algorithm 2: VLM-Verified Counterfactual Hindsight.}\\[2pt]
\textbf{Input.} Training env $E_{\text{train}}$; verifier env
$E_{\text{verify}}$ (separate \texttt{gym.make} instance, same env\_name);
VLM counterfactual function $\mathrm{CF}_\phi$; failure timestep $K$;
horizon $N\!=\!50$; success threshold $\eta\!=\!0.05$ m; reject-teleport
radius $\rho\!=\!0.05$ m; HER replay $\mathcal{D}_B$.\\[2pt]
\textbf{1. Trigger condition.} After a failed episode terminating at step
$T$ with $\|\text{ag}_T - g\| > \eta$.\\[2pt]
\textbf{2. Snapshot.} Capture
$s_K = (\text{qpos}_K, \text{qvel}_K, t_K, \text{mocap\_pos}_K, \text{mocap\_quat}_K, g)$
from $E_{\text{train}}$ at the localized failure timestep $K$ (from
$\mathrm{Loc}_\phi$ or oracle).\\[2pt]
\textbf{3. VLM counterfactual.} Query
$\text{out}\leftarrow \mathrm{CF}_\phi(\{F_k\}, K, g,
\text{task\_desc}, \text{variant}=\textsc{achieved\_goal})$ to receive
candidate corrective position $p_{\text{cf}}\in\mathbb{R}^3$ and (in the
\textsc{all} variant) corrective action $a_{\text{cf}}\in\mathbb{R}^4$.\\[2pt]
\textbf{4. Teleport gate.} If
$\|p_{\text{cf}}-g\|<\rho$: reject (reason=\texttt{teleport\_collapse});
fall back to standard HER and return.\\[2pt]
\textbf{5. Inverse-kinematic action (when needed).} If \textsc{achieved\_goal}
variant: synthesize $a_{\text{cf}}$ by a simple proportional controller
$a_{\text{cf}} = \mathrm{clip}((p_{\text{cf}}-\text{ee}_K)/0.05, -1, 1)$
appended with \texttt{grip}$=-1$ (close). If \textsc{action}/\textsc{all}:
use the VLM's $a_{\text{cf}}$ directly, clipped to $[-1,1]^4$.\\[2pt]
\textbf{6. Restore.} On $E_{\text{verify}}$: write back \texttt{qpos},
\texttt{qvel}, \texttt{time}, \texttt{mocap}; call
\texttt{mujoco.mj\_forward(model, data)} to propagate kinematic
derivatives. Verify round-trip equality $\|\text{qpos}_K -
\text{qpos}_K^{(\text{restored})}\|=0$ as a debug assertion.\\[2pt]
\textbf{7. Rollout.} For $j=0,\dots,N-1$: step
$(s', r_j)\leftarrow E_{\text{verify}}.\texttt{step}(a_{\text{cf}})$
under the env's \emph{unmodified} sparse reward
$r_j=E_{\text{verify}}.\texttt{compute\_reward}(\text{ag}_j, g)$.\\[2pt]
\textbf{8. Acceptance.} If $\exists j: r_j\ge 0$: VERIFIED. Set
$\text{ag}^\star\leftarrow\text{ag}_{j^\star}$. Append
$(s_K, a_K, \dots, s_{K+j^\star})$ to $\mathcal{D}_B$ with env's sparse
rewards \emph{and} a hindsight-relabeled copy with
$g'\leftarrow\text{ag}^\star$, both marked confidence$=$1.0. Else:
reject (\texttt{goal\_not\_reached}); fall back to standard HER.\\[2pt]
\textbf{Output.} (optional) verified trajectory $\tau_v$ appended to
buffer; verification flag $\in\{\text{ACCEPT},\text{REJECT}\}$.
}}
\caption{VLM-Verified Counterfactual Hindsight. The teleport gate
(step~4) is empirically necessary on planar tasks (FetchPush) but
structurally redundant on \textsc{action}-only variants where no
position field is emitted. The simulator-fork mechanism in step~6 has
zero modeling error (the verifier is the ground-truth env); per-call
CPU cost is 17.6\,ms.}
\label{alg:verified-cf}
\end{figure}

\subsection{Notation glossary}
\label{app:notation}

Table~\ref{tab:notation} gives the full notation used throughout the
paper. Each symbol's introduction section is noted in the rightmost
column.

\begin{table}[t]
\centering
\footnotesize
\caption{Complete notation glossary.}
\label{tab:notation}
\begin{tabular}{l p{0.62\linewidth} l}
\toprule
Symbol & Meaning & First use \\
\midrule
$\tau$, $T$ & Episode trajectory $\tau=(s_0,a_0,r_0,\dots,s_T)$; horizon $T\!=\!50$ for Fetch & \S\ref{sec:theory} \\
$i, \tau_i$ & Transition index in buffer; trajectory containing slot $i$ & \S\ref{sec:theory} \\
$s_t,a_t,r_t,g$ & State, action, reward, desired goal at step $t$ & \S\ref{sec:theory} \\
$\text{ag}_t$ & Achieved goal at step $t$ (object position for manipulation tasks) & \S\ref{sec:method:verified} \\
$\mathcal{D}_B$ & Replay buffer ($|\mathcal{D}_B|=10^6$ slots) & \S\ref{sec:theory} \\
$\theta,\pi_\theta,Q_\theta$ & SAC actor/critic parameters and policy/value heads & \S\ref{sec:exp:setup} \\
$\bar\theta$ & Target-critic parameters (Polyak average) & Alg.~\ref{alg:sem-per} \\
$\delta_i$ & TD-error at transition $i$ & Eq.~\eqref{eq:loss} \\
$\mu(i)$ & Sampling distribution over slots; subscripts $U,P,\text{Sem}$ & \S\ref{sec:theory} \\
$\alpha$ & PER priority exponent ($\alpha=0.6$) & Eq.~\eqref{eq:per-is} \\
$\beta_t$ & Annealed IS exponent ($\beta_0=0.4\to 1$ linear) & Eq.~\eqref{eq:per-is} \\
$\varepsilon$ & Priority floor (1e-6) & \S\ref{sec:method:sper} \\
$w_{\text{IS}}(i)$ & Per-sample IS-correction weight $(|\mathcal{D}_B|\mu(i))^{-\beta}$ & Eq.~\eqref{eq:per-is} \\
$t^\star,t_{\text{fail}}$ & Latent outcome-causing timestep; VLM's argmax thereof & \S\ref{sec:theory} \\
$q_\phi(t^\star\!\mid\!\tau)$ & VLM-emitted approximate posterior over outcome-causing timestep & Eq.~\eqref{eq:wsem} \\
$p(t^\star\!\mid\!\tau)$ & True (unknown) posterior over outcome-causing timestep & \S\ref{sec:theory} \\
$\phi$ & Frozen VLM parameters (e.g., GPT-4o, Claude Opus 4.7) & \S\ref{sec:theory} \\
$K_W(i;t^\star)$ & Window kernel; $1+(w_{\max}-1)\mathbb{1}[|i-t^\star|\le W]$ & Eq.~\eqref{eq:wsem} \\
$w_{\text{sem}}(i;\tau)$ & Per-transition semantic boost (Eq.~\ref{eq:wsem}) & Eq.~\eqref{eq:wsem} \\
$w_{\max}$ & Bounded boost cap ($w_{\max}=10$) & Eq.~\eqref{eq:wsem} \\
$W$ & Window half-width in timesteps ($W=5$) & Eq.~\eqref{eq:wsem} \\
$K$ & Number of keyframes per VLM query ($K=5$) & \S\ref{sec:method:sper} \\
$\mu_{\text{Sem}}(i)$ & Semantic-PER sampling distribution; product form (Eq.~\ref{eq:headline}) & Eq.~\eqref{eq:headline} \\
$\mathrm{CF}_\phi$ & VLM counterfactual function (4 prompt variants) & \S\ref{sec:method:prompt} \\
$p_{\text{cf}}, a_{\text{cf}}$ & VLM-proposed corrective position $\in\mathbb{R}^3$; corrective action $\in\mathbb{R}^4$ & \S\ref{sec:method:prompt} \\
$\rho$ & Reject-teleport radius (0.05\,m) & \S\ref{sec:method:prompt} \\
$\eta$ & Fetch goal-distance success threshold (0.05\,m) & \S\ref{sec:exp:setup} \\
$N$ & Verifier rollout horizon ($N\!=\!50$ steps, one Fetch horizon) & \S\ref{sec:method:verified} \\
$s_K$ & MuJoCo snapshot $(\text{qpos},\text{qvel},t,\text{mocap},g)$ at $K$ & \S\ref{sec:method:verified} \\
$\lambda_t$ & Sharony's mixture coefficient ($0\to 0.5$ warm-up) & \S\ref{sec:related} \\
$\gamma$ & MDP discount factor (0.98) & \S\ref{sec:exp:setup} \\
$\tau_{\text{SAC}}$ & Polyak target-update coefficient (0.005) & \S\ref{sec:exp:setup} \\
$k_{\text{HER}}$ & HER relabel multiplier ($k=4$, \texttt{future} strategy) & \S\ref{sec:exp:setup} \\
\bottomrule
\end{tabular}
\end{table}

\subsection{Heuristic localizer (``Oracle v3'') --- extended geometric reasoning}
\label{app:oracle-v3}

Section~\ref{app:prompts} (Appendix~A.2) of the main text introduces the
\texttt{GoalDistanceLocalizer} in three precedence-ordered phases. Here
we give the full per-environment geometric reasoning that drives those
phases.

\paragraph{FetchPush --- contact-loss detection.}
The Push task requires the gripper to maintain pushing contact with the
red cube along a contact normal aligned with the desired-goal direction.
We define the binary contact predicate
$\mathrm{ct}_t = \mathbb{1}[\|\text{ee}_t - \text{obj}_t\| < 0.05~\mathrm{m}]$
where $\text{ee}_t$ is the end-effector position and $\text{obj}_t$ the
object center. A contact-loss event at step $t$ is a transition
$(\mathrm{ct}_t,\mathrm{ct}_{t+1})=(1,0)$ such that
(i) the run length of $\neg \mathrm{ct}$ in $[t\!+\!1,\dots]$ is $\ge 2$
(filters single-frame contact bounces),
(ii) the object-goal distance $\|\text{obj}_t - g\| > 0.10$\,m
(filters near-success contact lifts),
(iii) the contact-loss timestep is at least 10 steps before the closest-approach
timestep (filters end-of-episode argmin coincidence).
On Push, contact loss frequently occurs because the agent
over-rotates the gripper or pushes past the cube perpendicular to the
target direction.

\paragraph{FetchPickAndPlace --- ballistic vs.\ argmin precedence.}
PickAndPlace has two distinct failure regimes: dropped-after-grasp
(ballistic-like) and missed-grasp (closest-approach). The localizer first
runs the ballistic-throw detector: compute per-step displacements
$d_t = \|\text{obj}_{t+1} - \text{obj}_t\|$, find peak $t_p\!=\!\arg\max d_t$
with $d_{t_p} > 0.008$\,m, and check that the average post-peak velocity
in the window $[t_p\!+\!1, t_p\!+\!7]$ exceeds $0.35\cdot d_{t_p}$. If
true, mark $t_{\text{fail}}\!=\!t_p$ (the drop moment). If false (object
remained ``carried'' or never moved), fall through to argmin: $\arg\min_t
\|\text{obj}_t - g\|$. This precedence is correct because a dropped
object's argmin-distance is dominated by the post-drop ballistic fall, which
is geometrically downstream of the causal release point.

\paragraph{FetchSlide --- release-point detection.}
Slide is the dominant ballistic environment: the puck slides across the
table after release, so contact is intentional and brief. The ballistic
detector almost always fires. We additionally require
(iv) the post-peak window mean exceeds the running median of pre-peak
displacements (filters spurious peaks from grasping motions), and
(v) the puck is in the table plane $z<0.45$\,m at peak time (filters
gripper-only motion). This identifies the release timestep, which on
Slide is the causal failure timestep: an off-axis strike at $t_p$ is
unrecoverable by step $T$.

\paragraph{Argmin fallback.}
If neither (b) nor (c) fires (the object stays still or moves slowly
relative to noise), we return $\arg\min_t \|\text{obj}_t - g\|$. This is
the closest-approach timestep, which on contact tasks where the agent
maintains pushing throughout is a reasonable diagnostic for ``the moment
divergence began''.

\subsection{Counterfactual prompt templates --- full text}
\label{app:cf-prompts}

We reproduce the four counterfactual prompt variants verbatim. All four
share the preamble (task description, $K$ keyframes, frame-to-timestep
map, failure-frame index) and differ only in the question / response
schema.

\paragraph{Common preamble.}
{\small
\begin{verbatim}
You are analyzing a robotic manipulation trajectory.

Task: {task_description}

{K} keyframes from a FAILED episode are shown in chronological order.
Frame-to-timestep mapping:
{frame_index_map}

A failure-detection module identified Frame {failure_frame_index}
(~{failure_frame_pct:.0f}% through the episode) as the critical moment
where the robot's action caused or revealed the failure.
\end{verbatim}}

\paragraph{Variant 1: \textsc{narrative}.}
{\small
\begin{verbatim}
QUESTION: At Frame {failure_frame_index}, what should the robot have
done differently for the episode to succeed? Be specific about the
corrective behaviour in physical terms (e.g. "lift higher before
approach", "close gripper sooner", "push from the opposite side").

Respond with ONLY a JSON object:
{
  "explanation": "<one concise sentence describing the corrective
                  behaviour>",
  "confidence": <float in [0,1]>
}
\end{verbatim}}

\paragraph{Variant 2: \textsc{action}.}
{\small
\begin{verbatim}
The robot's action space is a 4-D vector: (dx, dy, dz, gripper_command).
  - dx, dy, dz in [-1, 1] are end-effector velocity commands.
    Units: 1.0 ~= 5 cm/step. The 'x' axis points away from the robot,
    'y' to the left, 'z' upward.
  - gripper_command in [-1, 1]: -1 = close, +1 = open.

QUESTION: What 4-D corrective action vector should the robot have
executed at Frame {failure_frame_index} (instead of whatever it actually
did) so the episode would have succeeded? Provide a single delta-action
[dx, dy, dz, grip].

Be conservative: prefer small in-distribution deltas. If the failure is
a released-too-early issue, use grip ~= -1. If it's a missed-target issue,
point the (dx,dy,dz) toward the goal as seen in the frames.

Respond with ONLY a JSON object:
{
  "corrective_action": [<dx>, <dy>, <dz>, <grip>],
  "explanation": "<one concise sentence>",
  "confidence": <float in [0,1]>
}
\end{verbatim}}

\paragraph{Variant 3: \textsc{achieved\_goal}.}
{\small
\begin{verbatim}
In the Fetch suite, `achieved_goal` is the current 3D position of the
manipulated object (or the gripper, for FetchReach). The episode
succeeds when ||achieved_goal - desired_goal|| < 5 cm. Workspace coords
range roughly: x in [1.05, 1.55] m, y in [0.4, 1.1] m, z in [0.4, 0.85] m
(the table surface is at z ~= 0.42 m).

The desired_goal (target marker) for this episode is at:
  desired_goal = {desired_goal}

QUESTION: At Frame {failure_frame_index}, what 3D position SHOULD the
object have reached for the episode trajectory to be on track for
success? In other words, propose a hindsight `achieved_goal` for this
frame — a counterfactual location that, had the object been there at
this timestep, would have set the trajectory up to terminate within 5
cm of the desired_goal.

The proposed position must lie inside the Fetch workspace (above the
table).

Respond with ONLY a JSON object:
{
  "corrective_position": [<x>, <y>, <z>],
  "explanation": "<one concise sentence>",
  "confidence": <float in [0,1]>
}
\end{verbatim}}

\paragraph{Variant 4: \textsc{all} (unified).}
{\small
\begin{verbatim}
Context for the action and goal spaces:
  - Action: 4-D delta vector (dx, dy, dz, gripper), components in [-1,1].
    Each unit ~= 5 cm/step end-effector motion; gripper -1=close / +1=open.
  - achieved_goal / desired_goal: 3-D object (or gripper) positions in m.
    The workspace spans roughly x in [1.05,1.55], y in [0.4,1.1],
                                 z in [0.4,0.85].
  - desired_goal for this episode = {desired_goal}.

Provide a unified counterfactual:
  (a) the 3D position the object should have been at this frame
      (`corrective_position`),
  (b) the 4-D action the robot should have executed at this frame
      (`corrective_action`),
  (c) a one-sentence physical explanation.

Respond with ONLY a JSON object:
{
  "corrective_position": [<x>, <y>, <z>],
  "corrective_action": [<dx>, <dy>, <dz>, <grip>],
  "explanation": "<one concise sentence>",
  "confidence": <float in [0,1]>
}
\end{verbatim}}

\paragraph{Task description substitutions.}
The single \texttt{\{task\_description\}} placeholder takes one of three
strings:
\begin{itemize}\itemsep0pt
\item \texttt{FetchPush}: ``the robot should push the red cube to the green target on the table''.
\item \texttt{FetchPickAndPlace}: ``the robot should pick up the red cube and place it at the floating green target''.
\item \texttt{FetchSlide}: ``the robot should strike the red puck so that it slides to the green target''.
\end{itemize}

\section{Theoretical Derivations}
\label{app:theory}

\subsection{Importance-sampled posterior reweighting --- full derivation}
\label{app:is-derivation}

This section gives the step-by-step IS argument referenced in
Section~\ref{sec:theory}. The contribution of the framing is not the
derivation itself (which is standard self-normalized IS applied to a
prioritized-replay regression objective) but the identification of the
VLM as the proposal distribution and its location in the off-policy
correction taxonomy.

\paragraph{Step 1: the target objective.}
Off-policy value learning targets the Bellman-residual loss under
\emph{uniform replay}:
\begin{equation}
\mathcal{L}_U(\theta) \;=\; \mathbb{E}_{i\sim\mu_U}\bigl[\ell(\delta_i(\theta))\bigr]
\;=\; \frac{1}{|\mathcal{D}_B|}\sum_{i=1}^{|\mathcal{D}_B|} \ell(\delta_i(\theta)),
\label{eq:target}
\end{equation}
where $\ell$ is squared error or Huber. Equation~\eqref{eq:target} is
unbiased: the empirical-mean gradient targets the population gradient.

\paragraph{Step 2: importance-sampled estimator with a non-uniform proposal.}
For any proposal $\mu$ with $\mu(i)>0$ wherever $\mu_U(i)>0$, we have
the IS identity
\begin{equation}
\mathcal{L}_U(\theta) \;=\; \mathbb{E}_{i\sim\mu}\!\left[\frac{\mu_U(i)}{\mu(i)}\,
\ell(\delta_i(\theta))\right]
\;=\; \mathbb{E}_{i\sim\mu}\!\left[\frac{1}{|\mathcal{D}_B|\,\mu(i)}\,\ell(\delta_i(\theta))\right].
\label{eq:is-identity}
\end{equation}
This estimator is unbiased for $\mathcal{L}_U$ whenever the IS ratio is
non-zero on the support of $\mu_U$. The variance is minimized when $\mu$
is proportional to the integrand $\ell(\delta_i)$ --- precisely the
motivation for PER.

\paragraph{Step 3: PER's annealed estimator.}
Schaul et al.'s PER chooses
$\mu_P(i)\propto(|\delta_i|+\varepsilon)^\alpha$ and \emph{anneals}
the IS exponent:
\begin{equation}
\widehat{\mathcal{L}}_P(\theta) \;=\; \mathbb{E}_{i\sim\mu_P}\!\left[
\bigl(|\mathcal{D}_B|\,\mu_P(i)\bigr)^{-\beta_t}\,\ell(\delta_i(\theta))
\right], \qquad \beta_t\colon \beta_0=0.4 \to 1.
\label{eq:per-anneal}
\end{equation}
At $\beta_t=1$ this is exactly the unbiased estimator of
$\mathcal{L}_U$. For $\beta_t<1$ the estimator is biased (towards
emphasizing the high-priority tail) but lower-variance.

\paragraph{Step 4: Semantic PER as IS with a product proposal.}
We introduce the VLM posterior $q_\phi(t^\star\!\mid\!\tau)$ over
outcome-causing timesteps and the per-transition semantic boost
$w_{\text{sem}}(i;\tau_i)$ (Eq.~\ref{eq:wsem}). The Semantic-PER proposal
is the product
\begin{equation}
\mu_{\text{Sem}}(i) \;=\; \frac{\mu_P(i)\,w_{\text{sem}}(i;\tau_i)}
{\sum_j \mu_P(j)\,w_{\text{sem}}(j;\tau_j)}.
\label{eq:sem-prop}
\end{equation}
By Step~2, the strictly-correct IS-corrected estimator is
\begin{equation}
\widehat{\mathcal{L}}_{\text{Sem}}(\theta) = \mathbb{E}_{i\sim\mu_{\text{Sem}}}\!\left[
\bigl(|\mathcal{D}_B|\,\mu_{\text{Sem}}(i)\bigr)^{-\beta_t}\,\ell(\delta_i)\right].
\label{eq:sem-corrected}
\end{equation}
This estimator is the principled multiplicative reweighting referred to in
Equation~\eqref{eq:headline}. The factor $w_{\text{sem}}(i;\tau)$ enters
the proposal as a posterior-density modifier; the IS weight cancels it
asymptotically as $\beta_t\to 1$.

\paragraph{Step 5: the under-correction we actually run.}
In our current implementation we use the PER IS weight
$(|\mathcal{D}_B|\,\mu_P(i))^{-\beta_t}$ rather than the strictly-correct
$(|\mathcal{D}_B|\,\mu_{\text{Sem}}(i))^{-\beta_t}$. The ratio is
$(w_{\text{sem}}(i;\tau_i))^{\beta_t}\ge 1$: semantic-boosted slots are
under-corrected by exactly the boost factor (in $\beta_t$-power). The
resulting estimator is biased toward emphasizing the failure window.
This is mathematically analogous to V-trace's ratio truncation: a
controlled bias for variance reduction, with a known and computable
bias-amplification factor. The strictly-correct $w_{\text{IS},\text{Sem}}$
ablation has not been run (deferred to future work; cf.\ Limitations).

\subsection{Connection to V-trace and Retrace}
\label{app:vtrace}

V-trace \citep{espeholt2018impala} and Retrace($\lambda$)
\citep{munos2016retrace} are off-policy correction schemes for
\emph{multi-step temporal-difference returns} with the proposal
being the behavior policy. They specialize as follows.

\paragraph{Specialization of V-trace.}
V-trace defines per-step IS ratios $\rho_t = \pi(a_t\!\mid\!s_t)/\mu_b(a_t\!\mid\!s_t)$
and a target-policy product $c_t = \min(\bar c, \rho_t)$. The N-step
target is
$v_s = V(s) + \sum_{t\ge s}\gamma^{t-s}\bigl(\prod_{t'<t}c_{t'}\bigr) \min(\bar\rho,\rho_t)\,
\delta_t$. When $q_\phi$ is uniform on $\{0,\dots,T-1\}$, $w_{\text{sem}}\equiv 1$
and our scheme degenerates to PER --- which is single-step
($N\!=\!1$ in V-trace), with $\mu=\mu_P$ rather than $\mu_b$, and with
the IS weight $w_{\text{IS}}=(|\mathcal{D}_B|\mu_P)^{-\beta_t}$ playing
the role of V-trace's $\rho_t$. The truncation of V-trace
($\rho_t\to\min(\bar\rho,\rho_t)$) is conceptually identical to our
under-correction: a controlled bias for variance reduction.

\paragraph{Specialization of Retrace.}
Retrace's $c_t = \lambda\min(1,\rho_t)$ provides per-step trace cuts
controlled by $\lambda\in[0,1]$. Setting $\lambda=0$ reduces Retrace to
1-step Q-learning. In our notation, this is again the special case where
$q_\phi$ is uniform: the trace cut $c_t$ at every step is unity, no
re-weighting occurs, and the estimator reduces to TD(0) on the buffer
proposal.

\paragraph{Why the analogy is loose.}
V-trace and Retrace correct for \emph{policy mismatch} in multi-step
returns from a behavior trajectory. Semantic PER corrects for
\emph{replay-distribution mismatch} from a posterior-weighted proposal.
The proposals being corrected are different objects: V-trace's $\mu_b$
is the behavior policy at action selection; our $\mu_{\text{Sem}}$ is
the per-slot sampling distribution at minibatch construction. The
shared structure is the form of the correction --- IS ratios over a
proposal distribution, with controlled truncation --- not the
proposal's source.

\subsection{Conditions for unbiasedness}
\label{app:unbiased}

The semantic-PER estimator (Eq.~\ref{eq:sem-corrected}) is unbiased for
$\mathcal{L}_U$ when:

\begin{enumerate}\itemsep1pt
\item[(A1)] \textbf{Bounded boost.} $w_{\text{sem}}(i;\tau_i)\in[1,w_{\max}]$
with $w_{\max}<\infty$. This ensures $\mu_{\text{Sem}}\ll\mu_U$ (absolute
continuity) on the support of the buffer.
\item[(A2)] \textbf{Strict IS correction.} $\beta_t=1$ is used in
Equation~\eqref{eq:sem-corrected} (rather than the annealed value). For
$\beta_t<1$ the estimator is biased; the bias decays as $\beta_t\to 1$.
\item[(A3)] \textbf{Posterior independence of} $\theta$. $q_\phi$ does
\emph{not} depend on the current critic parameters $\theta$. In our
setting $q_\phi$ is conditioned on rendered keyframes from the
trajectory; the trajectory was collected under the policy at some prior
parameter value $\theta_{\text{collect}}\neq\theta_{\text{train}}$, but
$q_\phi$ itself is a frozen function of inputs. Assumption (A3) is
satisfied as long as the renderer does not depend on the critic.
\item[(A4)] \textbf{Posterior support.} $q_\phi(t^\star\!\mid\!\tau) > 0$
for every $t^\star\in\{0,\dots,T-1\}$ on which the VLM's argmax could
fall (in our window-based implementation, $q_\phi$ is automatically
non-zero everywhere because $w_{\text{sem}}\ge 1$ uniformly).
\end{enumerate}

The condition the framing does \emph{not} require is calibration of
$q_\phi$. Even an arbitrarily mis-calibrated VLM yields an unbiased
estimator of $\mathcal{L}_U$ provided (A1)--(A4) hold. Mis-calibration
affects \emph{variance}, not bias: an extremely concentrated $q_\phi$ at
a non-informative timestep would still give an unbiased estimator but
with very high variance (low effective sample size).

\paragraph{What our implementation gives up.}
In practice we use $\beta_t<1$ early in training (variance reduction)
and $w_{\text{IS},P}$ rather than $w_{\text{IS},\text{Sem}}$
(simpler implementation). This violates (A2) and yields a biased
estimator of $\mathcal{L}_U$. The bias points toward up-weighting the
failure window --- precisely the direction we want for credit
assignment. We frame this as a feature, not a bug, in line with
V-trace's truncation tolerance.

\section{Experimental Details}
\label{app:experiments}

\subsection{Hyperparameter tables}
\label{app:hparams}

Tables~\ref{tab:hparams-sac}--\ref{tab:hparams-vlm} list every
hyperparameter used in the paper. All values are constant across the
three Fetch environments unless explicitly stated; environment-specific
values are tabulated separately when needed.

\begin{table}[t]
\centering
\footnotesize
\caption{SAC hyperparameters (constant across envs).}
\label{tab:hparams-sac}
\begin{tabular}{l l l}
\toprule
Parameter & Value & Notes \\
\midrule
Hidden dim & 256 & Two-layer MLP for actor and critic \\
Learning rate (actor/critic/$\alpha$) & $3\!\times\!10^{-4}$ & Adam optimizer \\
Discount $\gamma$ & 0.98 & Fetch-standard \\
Polyak coefficient $\tau_{\text{SAC}}$ & 0.005 & Target-critic update \\
Initial temperature & 0.2 & Soft-actor entropy scale \\
Auto entropy tuning & enabled & Haarnoja et al.\ 2018 \\
Batch size & 256 & \\
Updates per env step & 1 & \\
Total env steps & 500k--1M & Paper-grade sweeps; 50k--100k for in-flight \\
Warm-up steps & 10k & Random-policy data before SAC update \\
Eval interval & 5k steps & 20 eval episodes, deterministic policy \\
Random seeds & 42, 123, 999 & 3 seeds per (method, env) cell \\
Replay capacity & $10^6$ & Ring buffer with sum-tree priorities \\
\bottomrule
\end{tabular}
\end{table}

\begin{table}[t]
\centering
\footnotesize
\caption{HER hyperparameters.}
\label{tab:hparams-her}
\begin{tabular}{l l l}
\toprule
Parameter & Value & Notes \\
\midrule
Strategy & \texttt{future} & Andrychowicz et al.\ 2017 \\
Relabel multiplier $k_{\text{HER}}$ & 4 & 4 hindsight transitions per real transition \\
Goal sampler & uniform over future ts & \\
\bottomrule
\end{tabular}
\end{table}

\begin{table}[t]
\centering
\footnotesize
\caption{Prioritized Replay (PER) hyperparameters.}
\label{tab:hparams-per}
\begin{tabular}{l l l}
\toprule
Parameter & Value & Notes \\
\midrule
$\alpha$ (priority exponent) & 0.6 & Schaul et al.\ 2016 \\
$\beta_0$ (initial IS exponent) & 0.4 & Linear anneal \\
$\beta_{\text{end}}$ & 1.0 & Hit at \texttt{per\_beta\_anneal\_steps} \\
$\beta$ anneal length & 500{,}000 steps & \\
$\varepsilon$ (priority floor) & $10^{-6}$ & \\
\bottomrule
\end{tabular}
\end{table}

\begin{table}[t]
\centering
\footnotesize
\caption{Semantic PER and Verified-CF hyperparameters.}
\label{tab:hparams-semper}
\begin{tabular}{l l l}
\toprule
Parameter & Value & Notes \\
\midrule
Window half-width $W$ & 5 timesteps & Boost around $t_{\text{fail}}$ \\
Boost cap $w_{\max}$ & 10 & Multiplicative on PER priority \\
Keyframes $K$ & 5 & Uniformly sampled per failed episode \\
VLM call interval & every failed episode & \texttt{vlm\_call\_interval}=1 \\
VLM provider & GPT-4o (production) & \texttt{gpt-4o-2024-08-06} \\
Heuristic localizer & GoalDistanceLocalizer & Privileged-state oracle \\
\midrule
\multicolumn{3}{l}{\textit{Verified-CF}} \\
Verifier rollout horizon $N$ & 50 steps & One Fetch episode \\
Verifier success threshold $\eta$ & 0.05 m & Fetch sparse-reward threshold \\
Teleport reject radius $\rho$ & 0.05 m & $\|p_{\text{cf}}-g\| < \rho \Rightarrow$ reject \\
CF prompt variant (production) & \textsc{achieved\_goal} & cf.\ Sec.~\ref{sec:exp:prompt} \\
Min VLM self-confidence & 0.5 & Soft filter; under-utilized \\
Verifier env action repeat & 1 (default) & \\
Verifier env reseed strategy & disabled & State restored explicitly \\
\bottomrule
\end{tabular}
\end{table}

\begin{table}[t]
\centering
\footnotesize
\caption{VLM call parameters.}
\label{tab:hparams-vlm}
\begin{tabular}{l l l}
\toprule
Parameter & Value & Notes \\
\midrule
Max output tokens & 512 & Sufficient for JSON + reasoning \\
Temperature (Claude) & not set & Opus 4.7 deprecates \texttt{temperature} kwarg \\
Temperature (GPT-4o) & 0.0 & Greedy decoding \\
Image detail (GPT-4o) & low & 512$\times$512 JPEG-encoded keyframes \\
Image quality (JPEG) & 85 & Encoder setting \\
Retry policy & 3 retries, exponential backoff & \\
Cost per call & $\le\!\$0.05$ & \texttt{gpt-4o-2024-08-06}; Opus 4.7 similar \\
Median wall-clock per call & 11s (Opus) / 4s (GPT-4o) & \\
\bottomrule
\end{tabular}
\end{table}

\subsection{Compute used}
\label{app:compute}

\paragraph{Hardware.}
Two compute substrates were used. \emph{Local}: a single RTX 5070~Ti
(16~GB) workstation with 64~GB DDR5 RAM and a 16-core Zen~5 CPU,
running Ubuntu 24.04. \emph{Cloud}: Modal A10G (24~GB) instances with
PyTorch~2.4 + CUDA~12.4. The headline sweeps in
Figure~\ref{fig:headline} are mixed local + Modal; the in-flight HER
baselines (27 SAC runs) are all Modal.

\paragraph{Wall-clock budget.}
A single 500k-step SAC+HER+Semantic-PER run on FetchPickAndPlace takes
$\approx 3.5$~h on A10G, end-to-end including VLM calls. The dominant
costs are (i) the SAC critic update ($\approx 60\%$), (ii) MuJoCo
stepping ($\approx 25\%$), (iii) VLM API calls (asynchronous,
$\approx 10\%$ when in-flight; budget-bound at $\$0.05\!\times\!1700$
failed eps $\approx \$85$/run); (iv) per-verification CPU
($\approx 5\%$, 17.6\,ms per verification). The 27-run HER sweep is
projected to total $\approx 100$~GPU-hours.

\paragraph{Total compute used (paper-grade).}
Approximate cumulative across all reported runs:
\begin{itemize}\itemsep0pt
\item \textbf{Local GPU}: $\approx 80$~h on RTX 5070~Ti.
\item \textbf{Modal A10G GPU}: $\approx 130$~h (HER baselines + in-flight verified-CF).
\item \textbf{CPU}: $\approx 200$~hours total (rendering, evaluation, analysis scripts).
\item \textbf{VLM API}: $\approx \$320$ across all prompt-design,
cross-task, real-data, and in-flight experiments. Largest single line:
in-flight verified-CF runs at $\approx \$85$/run.
\end{itemize}

\subsection{Random seeds and statistical methodology}
\label{app:stats}

\paragraph{What counts as a ``seed''.}
A seed is a single non-negative integer used to instantiate (i)
\texttt{numpy.random.default\_rng}, (ii)
\texttt{torch.manual\_seed}, (iii) \texttt{torch.cuda.manual\_seed\_all},
(iv) Gymnasium-Robotics \texttt{env.reset(seed=...)}, and (v) any
\texttt{random.Random} instances used by the trainer. SAC parameter
initialization is seeded by (iii); the replay buffer's slot-overwrite
order is seeded by (i); the env's per-episode reset (object position,
goal position) is seeded by (iv). Three seeds (42, 123, 999) are run per
(method, env) cell unless noted otherwise.

\paragraph{Confidence intervals.}
The reported error bars throughout the paper are mean~$\pm$~standard
error of the mean (SE = $s/\sqrt{n}$ where $s$ is the sample standard
deviation). For $n=3$ this is wide and only directional; we say so in
the figure captions. The cross-task transfer pilot
(Section~\ref{sec:exp:transfer}) reports raw fractions at $n=12$; a
12/12 binary fraction is not statistically distinguishable from an 80\%
underlying rate without a larger sample. Prompt-design metrics in
Table~\ref{tab:prompt-variants} are pooled across $\le\!6$ episodes per
cell; we report SE rather than bootstrap CIs because the per-episode
spread dominates and a bootstrap at $n=6$ underestimates uncertainty.

\paragraph{Bootstrap protocol (when used).}
Where bootstrap is invoked (figures with mean$\pm$band shading), we use
the basic non-parametric bootstrap with $B=1{,}000$ resamples drawn with
replacement at the seed level (i.e., resampling whole training-curves
rather than individual eval-episode successes). The reported band is
the 16th--84th percentile range, corresponding approximately to
$\pm 1\,\mathrm{SE}$ under a normal approximation.

\subsection{Evaluation protocol}
\label{app:eval}

\paragraph{Eval schedule.}
Eval is triggered every $5{,}000$ env steps (\texttt{eval\_interval}).
Each eval comprises $E\!=\!20$ episodes with the policy in deterministic
mode (action $=$ tanh of actor's mean output, no entropy bonus). The
env is reset with a fresh seed drawn from a stream independent of the
training seed.

\paragraph{Success criterion.}
An eval episode is a success iff $\|\text{ag}_T - g\|<0.05$~m at the
final step. The eval-time success rate is the fraction of the 20
episodes that succeeded. Note that Fetch's training reward is
$r=-\mathbb{1}[\|\text{ag}_t-g\|>0.05]$; eval success rate is reported
at $t=T$ only, not as the per-step reward, so a transient close-pass is
not credited.

\paragraph{Checkpoint frequency.}
Model checkpoints are saved every $50{,}000$ steps for offline analysis
and reproduction. Eval is logged independently (no model save at eval
time).

\paragraph{Final reporting metric.}
The reported headline metric is success rate at the \emph{final
checkpoint} (i.e., the eval-time success at the end of training for
each run). The plotted curves in Figure~\ref{fig:headline} show the
training-time trajectory of this metric, mean$\pm$SE across the three
seeds, with horizontal axis env-steps.

\section{Additional Results}
\label{app:results}

\subsection{Learning curves over training}
\label{app:learning-curves}

The headline curves in Figure~\ref{fig:headline} compare four conditions
(Uniform, PER, Semantic-PER (heuristic Oracle), Semantic-PER (GPT-4o))
on the three Fetch environments. In this appendix we note three
additional cuts that complement the figure.

\paragraph{Per-env final success rates.}
Table~\ref{tab:final-success} summarizes the eval-time success
rate at the final checkpoint, per (method, env, seed) cell.
Sample size is $n=3$ seeds; we report individual seed values rather
than aggregated SEs to give reviewers the underlying data. The Oracle
column is the heuristic privileged-state localizer; the GPT-4o column
uses the pre-fix \textsc{all} prompt variant (cf.\ pipeline-bug
discussion in Section~\ref{sec:discussion}).

\begin{table}[t]
\centering
\footnotesize
\caption{Per-seed final eval-time success rates (\%). Pre-fix
pipeline-bug numbers retained for transparency; the corrected
verified-CF curves are in flight and will replace the GPT-4o column in
the next revision.}
\label{tab:final-success}
\begin{tabular}{ll cccc}
\toprule
Env & Seed & Uniform & PER & Sem-PER (Oracle) & Sem-PER (GPT-4o, pre-fix) \\
\midrule
\multirow{3}{*}{FetchPush}        & 42  & 42 & 51 & 67 & 56 \\
                                  & 123 & 38 & 49 & 70 & 49 \\
                                  & 999 & 41 & 54 & 65 & 51 \\
\multirow{3}{*}{FetchPickAndPlace}& 42  & 18 & 33 & 51 & 47 \\
                                  & 123 & 21 & 31 & 53 & 49 \\
                                  & 999 & 16 & 35 & 54 & 50 \\
\multirow{3}{*}{FetchSlide}       & 42  & 11 & 18 & 22 & 14 \\
                                  & 123 & 9  & 15 & 21 & 17 \\
                                  & 999 & 13 & 20 & 24 & 12 \\
\bottomrule
\end{tabular}
\end{table}

\paragraph{Ablation placeholder.}
Section~\ref{sec:discussion} flags a pipeline-bug discovery that
re-frames the GPT-4o column in Table~\ref{tab:final-success}. The
corrected curves (verified-CF with the \textsc{achieved\_goal} variant
and 5\,cm gate enforced) are running overnight in the Modal cluster;
the corresponding entries are to be filled in the morning by the
data-collection agent. Pre-registered prediction: the GPT-4o-corrected
curve tracks the Oracle within 5--8\,pp of final success rate on
PickAndPlace, with similar tightening on Push and Slide.

\subsection{Cross-task transfer details}
\label{app:transfer}

Section~\ref{sec:exp:transfer} reports 12/12 task-relevant annotations
across 4 episodes/env on FetchPush, FetchPickAndPlace, and FetchSlide.
The judge's per-episode rationale across the three envs:
\begin{itemize}\itemsep1pt
\item Push: ``gripper passes over block without making contact'' (3/4),
``pushed block off the table edge'' (1/4).
\item PickAndPlace: ``gripper at block location but fails to close''
(2/4), ``block fell during transport'' (1/4), ``approach trajectory
misses block'' (1/4).
\item Slide: ``strike direction off-axis'' (2/4), ``strike velocity
insufficient for goal distance'' (1/4), ``puck strikes wall and
deflects'' (1/4).
\end{itemize}
Tolerant keyframe agreement (VLM keyframe argmax within $\pm 1$ of the
heuristic-Oracle keyframe argmax) improves monotonically with
visual-distinctiveness (Push 2/4, PickAndPlace 3/4, Slide 4/4),
suggesting the heuristic Oracle under-localizes synthetic-policy
rollouts on contact-heavy tasks rather than the VLM erring on visually
ambiguous frames.

\section{Reproducibility Checklist}
\label{app:repro}

\subsection{NeurIPS 2024 reproducibility checklist}

The NeurIPS 2024 reproducibility checklist asks for explicit answers
to a structured set of items. We answer each in turn.

\paragraph{Models, datasets, and benchmarks.}
\begin{itemize}\itemsep0pt
\item \textbf{Datasets used.} Gymnasium-Robotics \texttt{FetchPush-v4},
\texttt{FetchPickAndPlace-v4}, \texttt{FetchSlide-v4} (simulation only;
no human data). No new datasets are introduced.
\item \textbf{License.} Gymnasium-Robotics is MIT-licensed; the
underlying MuJoCo dependency is Apache-2.0. We comply with the public
license terms.
\item \textbf{Models used.} SAC trained from scratch; frozen VLMs
GPT-4o (\texttt{gpt-4o-2024-08-06}) and Claude Opus 4.7
(\texttt{claude-opus-4-7-20250529}). No new model weights are released.
\item \textbf{Train/eval split.} N/A. RL setting; eval is on held-out
seeds drawn from the same env distribution.
\end{itemize}

\paragraph{Algorithmic claims.}
\begin{itemize}\itemsep0pt
\item \textbf{Theoretical claims justified.} Section~\ref{sec:theory}
gives the IS-posterior framing; Appendix~\ref{app:is-derivation} gives
the step-by-step derivation. Assumptions (A1)--(A4) for unbiasedness
are stated in Appendix~\ref{app:unbiased}.
\item \textbf{Algorithm pseudocode.} Algorithms~\ref{alg:sem-per} and
\ref{alg:verified-cf} give full step-by-step pseudocode.
\item \textbf{Hyperparameters.} Tables~\ref{tab:hparams-sac}--\ref{tab:hparams-vlm}
list every hyperparameter used.
\end{itemize}

\paragraph{Code and data.}
\begin{itemize}\itemsep0pt
\item \textbf{Code release.} The codebase is released at
\url{https://github.com/d-grant-uc-berkeley/RL_project}, commit
\texttt{0fa36fc}, with subdirectories \texttt{src/} (libraries),
\texttt{configs/} (YAML training configs), \texttt{scripts/}
(experiment harnesses), and \texttt{agent\_reports/} (analysis
artifacts). The verified-CF prototype is at
\texttt{src/vlm/verified\_counterfactual.py}; the smoke harness is at
\texttt{scripts/n1\_smoke\_verified\_cf.py}.
\item \textbf{Data release.} All counterfactual-prompt outputs
(\texttt{C1\_counterfactual\_outputs.json},
\texttt{C1v2\_gpt4o\_outputs.json},
\texttt{C1v2\_real\_data\_outputs.json}, and the cross-task validation
outputs) are released in the same repository under
\texttt{agent\_reports/}.
\item \textbf{Anonymization.} For double-blind review the GitHub URL
above will be replaced by an anonymized Zenodo bundle. The full
training logs (W\&B project ``RL\_project'' under entity
\texttt{d-grant-uc-berkeley}) are private during review; they will be
made public on acceptance.
\end{itemize}

\paragraph{Compute.}
\begin{itemize}\itemsep0pt
\item Local: RTX 5070~Ti (16~GB), 16-core Zen~5 CPU, 64~GB RAM, Ubuntu
24.04. Cloud: Modal A10G (24~GB) instances.
\item Software: Python~3.10, PyTorch~2.4 + CUDA~12.4,
Gymnasium~1.0.0, Gymnasium-Robotics~1.3.0, MuJoCo~3.1.6,
NumPy~1.26, Anthropic SDK 0.39, OpenAI SDK 1.55.
\item Total compute: $\approx 80$~h local GPU, $\approx 130$~h Modal
A10G, $\approx 200$~h CPU, $\approx \$320$ in VLM API calls. See
Appendix~\ref{app:compute} for breakdown.
\end{itemize}

\paragraph{Statistical reporting.}
\begin{itemize}\itemsep0pt
\item Sample sizes are explicitly noted ($n\!=\!3$ seeds; $n\!=\!4$--6 per
prompt-variant cell; $n\!=\!12$ cross-task pilot). All small-sample
claims are flagged as preliminary in Section~\ref{sec:discussion}.
\item Confidence intervals: mean~$\pm$~SE throughout, with bootstrap
where seed-level resampling is feasible (Appendix~\ref{app:stats}).
\end{itemize}

\paragraph{Ethics.}
\begin{itemize}\itemsep0pt
\item No human subjects, no personally identifying data, no demographic
attributes.
\item VLM API usage is short, low-volume, and unrelated to any
personal-data domain.
\item Simulation only; no physical robot deployment.
\end{itemize}

\subsection{Broader Impacts}
\label{app:impact}

Sparse-reward robotic manipulation is a methodological building block
of robot learning. The mechanisms introduced here --- Semantic PER and
VLM-Verified Counterfactual Hindsight --- improve sample efficiency on
simulated manipulation tasks without changing the policy class,
deployment surface, or safety properties of the underlying RL agent.

\paragraph{Positive impacts.}
(i) Reduced sample requirements lower the energy footprint of robot
learning research, both in simulation and (downstream) in
sim-to-real settings. (ii) The IS-posterior framing provides a
principled lens for VLM-in-replay methods, which we hope inspires more
mathematically grounded subsequent work. (iii) The verified-CF
acceptance gate is a general-purpose pattern (``language prior +
physics test'') that may transfer to other domains where foundation
models output low-bandwidth structured suggestions.

\paragraph{Negative impacts and mitigations.}
(i) Foundation-model API use carries a non-trivial energy and dollar
cost; we mitigate by reporting per-run API spend
(Appendix~\ref{app:compute}) and providing the heuristic-Oracle
fallback that requires no API calls. (ii) The VLM may hallucinate
failure modes that are physically implausible; the verification gate
(Algorithm~\ref{alg:verified-cf}) is our primary mitigation. (iii) No
robot-deployment safety risks are introduced beyond standard
sim-trained policy caveats.

\paragraph{Scope.}
This work is methodological and simulation-only. We make no claims
about real-robot performance, sim-to-real transfer, or any specific
deployment scenario.

\subsection{Compute infrastructure details}
\label{app:infra}

For full reproducibility, the exact environment specifications used:
\begin{itemize}\itemsep0pt
\item \textbf{OS}: Ubuntu 24.04 LTS (kernel 6.17), glibc 2.39.
\item \textbf{Python}: 3.10.14, managed via the project's \texttt{.venv} (pip).
\item \textbf{PyTorch}: 2.4.1+cu124 on CUDA 12.4, cuDNN 9.1.
\item \textbf{Gymnasium}: 1.0.0; \textbf{Gymnasium-Robotics}: 1.3.0;
\textbf{MuJoCo}: 3.1.6; \textbf{mujoco-py}: not used (pure MuJoCo Python bindings).
\item \textbf{Rendering}: EGL backend (\texttt{MUJOCO\_GL=egl}) for
headless Modal containers; OSMesa for local debugging.
\item \textbf{VLM SDKs}: \texttt{anthropic} 0.39.0,
\texttt{openai} 1.55.0; both pinned in
\texttt{requirements.txt}.
\item \textbf{Other}: \texttt{numpy} 1.26.4, \texttt{Pillow} 11.0.0,
\texttt{pyyaml} 6.0.2, \texttt{wandb} 0.18.7.
\item \textbf{Modal image}: \texttt{python:3.10-slim} base, plus
\texttt{libegl1-mesa libgl1 libosmesa6 libosmesa6-dev}.
\end{itemize}

A \texttt{Dockerfile} reproducing this exact environment is included in
the codebase release under \texttt{docker/} (for non-Modal users).

\section{Failure Mode Analysis}
\label{app:failures}

\subsection{The teleport-collapse failure mode (C1v2-A)}
\label{app:teleport}

\paragraph{Formal characterization.}
The teleport-collapse failure mode of VLM counterfactual prompting is
the event $\|p_{\text{cf}}-g\|<\rho$ where $p_{\text{cf}}$ is the VLM's
emitted \texttt{corrective\_position} for a failed episode and $g$ is
the desired-goal of that episode (with $\rho\!=\!0.05$~m). Empirically,
this mode appears in:
\begin{itemize}\itemsep0pt
\item Claude Opus 4.7, FetchPush, \textsc{achieved\_goal}: $4/4$ episodes (100\%).
\item Claude Opus 4.7, FetchPickAndPlace, \textsc{achieved\_goal}: $3/4$ episodes on real, $2/2$ on synthetic.
\item GPT-4o, FetchPush, \textsc{achieved\_goal}: $0/6$ on synthetic;
on the joint \textsc{all} variant: $4/6$ on synthetic.
\item Claude Sonnet 4.5, FetchPush, \textsc{achieved\_goal}: $3/4$ on real (75\%).
\end{itemize}

\paragraph{Why it happens.}
The \textsc{achieved\_goal} prompt asks for ``the position the object
should have reached at the failure frame for the episode to be on
track''. When the desired goal is on the table surface (FetchPush) and
the failure frame is mid-episode, there is no non-trivial waypoint:
``on track'' is structurally equivalent to ``at the target now''. The
VLM resolves this ambiguity by output $p_{\text{cf}}=g$ exactly. On
PickAndPlace, the goal is in mid-air ($z>0.43$~m), so a non-trivial
waypoint exists (the gripper-and-block at mid-height) and the
teleport mode is partially evaded.

\paragraph{Mitigations evaluated.}
\begin{enumerate}\itemsep0pt
\item \textbf{Hard rejection gate.} Reject all outputs with
$\|p_{\text{cf}}-g\|<\rho$. Effective at gate-time but loses signal
($\sim$100\% of FetchPush \textsc{achieved\_goal} outputs rejected for
Claude Opus 4.7). Adopted as a belt-and-suspenders fallback in the
production pipeline.
\item \textbf{Prompt-variant switch to \textsc{all}.} The unified
\textsc{all} prompt structurally regularizes the position field (must
co-emit a position \emph{and} a 4-D action), reducing teleport rate
on real data from 100\% to 17--25\%. Helpful but not eliminative.
\item \textbf{Model switch.} Switching Opus~4.7 to Sonnet~4.5
rescues PickAndPlace (75\% $\to$ 0\%) but not Push (75\% persists).
Switching to GPT-4o eliminates Push on synthetic (100\% $\to$ 0\%)
but persists on the joint \textsc{all} variant.
\item \textbf{Mechanism replacement (verified-CF).} Switch to the
\textsc{action} variant + simulator-fork verification
(Algorithm~\ref{alg:verified-cf}). Teleports are structurally
inexpressible: a 4-vector cannot teleport a 50~g block by 50~cm. This is
the adopted mechanism.
\end{enumerate}

\paragraph{Adopted production configuration.}
\texttt{(GPT-4o, \textsc{achieved\_goal} with $\rho\!=\!0.05$ gate,
verified-CF fallback)}. The combination of (i) GPT-4o's strict
dominance on the synthetic \textsc{achieved\_goal} grid, (ii) the gate
catching residual teleports, and (iii) the verified-CF mechanism
absorbing rejections gives a defense-in-depth strategy.

\subsection{Action sign-flip in real-data evaluation (C1v2-B)}
\label{app:sign-flip}

\paragraph{Phenomenon.}
On the \textsc{action} prompt variant, $\sim\!50$\% of VLM-emitted
4-vectors have at least one axis whose sign disagrees with
$\mathrm{sign}(g - \text{ag}_K)$ for the failure-frame achieved goal.
The VLM's textual explanation \emph{correctly} describes the failure
mode in $4/4$ episodes (e.g., ``descend and close the gripper''), but
the emitted action vector points in the wrong axis-direction. The judge
catches each contradiction with phrasing like ``the action's negative-y
component moves away from the target''.

\paragraph{When it occurs.}
The sign-flip rate per (source, model, variant):

\begin{itemize}\itemsep0pt
\item Synthetic, Opus 4.7, \textsc{action}: 0.50 ($n\!=\!4$).
\item Synthetic, Opus 4.7, \textsc{all}: 0.17 ($n\!=\!2$).
\item Real, Opus 4.7, \textsc{action}: 0.55 ($n\!=\!10$).
\item Real, Sonnet 4.5, \textsc{action}: 0.70 ($n\!=\!5$).
\item Real, Opus 4.7, \textsc{all}: 0.39 ($n\!=\!9$).
\end{itemize}

\paragraph{Hypothesized cause.}
VLMs have known difficulty grounding metric 3-D axis directions
(``which way is +y'') from 2-D rendered keyframes. The Fetch envs use
an angled tabletop camera; the y-axis appears as a horizontal
in-plane axis with no visual anchor, and the VLM frequently confuses
its sign. The \textsc{all} variant reduces this rate by forcing the VLM
to co-emit a position field (which it tends to get right) alongside
the action, providing implicit cross-checking.

\paragraph{Recommended sign-check gate.}
The recommended mitigation for any pipeline consuming \textsc{action}
outputs is the simple gate
\begin{equation}
\text{accept iff } \mathrm{sign}(a_{\text{cf}}[:\!3]) = \mathrm{sign}(g - \text{ag}_K)
\quad \text{component-wise on axes with } |g - \text{ag}_K| > 0.01~\mathrm{m}.
\label{eq:sign-gate}
\end{equation}
Empirically this filters $\approx 50$\% of \textsc{action} outputs.
The remaining $\sim\!50$\% have correct signs and are forwarded to
either the verified-CF mechanism or direct action-relabeling.

\paragraph{Why the verified-CF mechanism subsumes this.}
The simulator-fork verification (Algorithm~\ref{alg:verified-cf}) is
\emph{strictly stronger} than the sign-check gate: an
$a_{\text{cf}}$ with a flipped sign on any axis points the gripper
away from the goal, so the verifier env rolls out for 50~steps without
firing the sparse reward, and the counterfactual is rejected as
\texttt{goal\_not\_reached}. The sign-check gate is therefore mostly a
diagnostic and a pre-filter to save the verifier's 17.6\,ms per call;
on the production pipeline the verifier alone suffices.

\section{Reproducibility commands}
\label{app:commands}

For users wishing to reproduce specific results:

\paragraph{Headline ablation (Figure~\ref{fig:headline}).}
{\small
\begin{verbatim}
# Local single-seed smoke
python train.py --config configs/her_semantic_per.yaml \
    --env.name FetchPickAndPlace-v4 --training.seed 42 \
    --training.total_steps 50000

# Full Modal sweep (27 runs)
modal run modal_app.py::run_ablation
\end{verbatim}}

\paragraph{Counterfactual prompt evaluation
(Tables~\ref{tab:prompt-variants}).}
{\small
\begin{verbatim}
# Synthetic (6 episodes)
python scripts/test_counterfactual.py \
    --provider openai --model gpt-4o \
    --judge_model claude-opus-4-7 \
    --variants narrative action achieved_goal all \
    --n_episodes 6 --budget 60

# Real (10 episodes)
python scripts/run_c1v2_real_data_eval.py \
    --n_episodes 10 --providers openai anthropic
\end{verbatim}}

\paragraph{Cross-task transfer (Figure~\ref{fig:transfer}).}
{\small
\begin{verbatim}
python scripts/n2_validate_cross_task.py \
    --envs FetchPush-v4 FetchPickAndPlace-v4 FetchSlide-v4 \
    --n_per_env 4 --K 5
\end{verbatim}}

\paragraph{Verified-CF smoke (Figure~\ref{fig:protocol}).}
{\small
\begin{verbatim}
python scripts/n1_smoke_verified_cf.py
# 4/4 expected pass; results in agent_reports/N1_smoke_results.json
\end{verbatim}}
